\section{Preliminaries}

\subsection{Grassmannians}
We denote the complex Grassmannian of sub-vector spaces of $\CC^n$ of dimension $r$ by $G(r,n)$. It is a smooth projective variety embedded into $\PP^{\binom{n}{r}-1}$ by the well known Plücker embedding; see \cite{EisHar}. For a point \(x\in G(r,n)\) we denote the corresponding \(r\)-dimensional subspace of \(\CC^n\) by \(L(x)\). The algebraic torus $T = (\CC^*)^n$ acts on $G(r,n)$ by dilating the coordinates of $\CC^n$. We will be interested in the closure in $G(r,n)$ of the $T$ orbit of $x \in G(r,n)$ denoted $\T$.

\subsection{Partitions}
\begin{definition}
    A partition $\lambda$ is a tuple $\lambda = (\lambda_1, \dots , \lambda_k)$ of positive integers with $\lambda_1 \geq \lambda_2 \geq \dots \geq \lambda_k$.
\end{definition}
We think of partitions in terms of their Young diagrams (sometimes called Ferrers diagrams), i.e.\ a collection of boxes with $\lambda_i$ boxes in row $i$. We write $|\lambda| = \sum_i\lambda_i$ for the number of boxes. There is a partial ordering of partitions defined by $\lambda \subseteq \mu$ if the Young diagram of $\lambda$ is contained in the Young diagram of $\mu$. We will be interested in the partitions $\lambda \subseteq r\times(n-r)$ where $r\times (n-r)$ is the partition whose Young diagram is the $r \times (n-r)$ rectangle. For each $\lambda \subseteq r\times (n-r)$ let $\lambda^c$ be the partition given by the complement of the Young diagram of $\lambda$ in the $r \times (n-r)$ rectangle, rotated by 180 degrees. We denote the so called hook partition by $h$, i.e.\ $h = (n-r,1,\dots,1)$ with $r-1$ ones. Notice that $h^c = (r-1)\times(n-r-1)$. The Young diagrams of $h$ and $h^c$ for $r=3$ and $n=7$ are shown below.
\begin{equation*}
        \scalebox{0.8}{\young(~~~~,~,~)} \hspace{2cm} \scalebox{0.8}{\young(~~~,~~~)}
    \end{equation*}
\begin{definition}
    A Young tableau of shape $\lambda$ is a filling of the boxes of the Young diagram of $\lambda$ with the integers $1,2,\dots, |\lambda|$. A Young tableau is called standard if it is increasing in rows and columns.
\end{definition}
The combinatorics of partitions and Young tableaux appears in many areas of mathematics, such as in the description of the Chow ring of Grassmannians.

\subsection{Chow ring of Grassmannians}
For a partition \(\lambda \subseteq r\times (n-r)\) the associated jumping sequence \(J = (j_i)_{i = 1}^r\) is defined by \(j_i = n-r+i-\lambda_i\). To obtain \(J\) from \(\lambda\), label the steps on the path from \((n-r,r)\) to \((0,0)\) along the Young diagram of \(\lambda\) with the numbers \(1,2,\dots , n\) in order. Then \(j_i\) is the label of the \(i\)th vertical step. The Schubert variety associated to \(\lambda\) and a complete flag \(V_\bullet\) of \(\CC^n\) is defined as 
\[
    X_\lambda(V_\bullet) = \left\{ x \in G(r,n) \mid \dim(L(x) \cap V_{j_i}) \geq i \right\}.
\]
Notice that if \(j_{i+1} = j_{i} + 1\) then
\[
    \dim(L(x) \cap V_{j_{i+1}}) \geq i+1 \implies \dim(L(x) \cap V_{j_i}) \geq i.
\]
So in this situation we may remove the \(i\)th defining inequality of \(X_\lambda(V_\bullet)\). The Chow ring $A^\bullet(G(r,n))$ is generated as a free abelian group by the Schubert cycles $\sigma_\lambda := [X_\lambda(V_\bullet)]$ where $\lambda$ ranges over all partitions $\lambda \subseteq r\times(n-r)$. Here $\sigma_\lambda$ is in the degree $|\lambda|$ piece of $A^\bullet(G(r,n))$ i.e.\ $\sigma_\lambda \in A^{|\lambda|}(G(r,n))$. The top piece $A^{r(n-r)}(G(r,n))$ is generated by the Schubert cycle $\sigma_{r\times(n-r)}$ and the isomorphism with $\Z$ sending $\sigma_{r\times(n-r)}$ to $1$ is called the degree map $\deg\colon A^{r(n-r)}(G(r,n)) \to \Z$.

The product in $A^\bullet(G(r,n))$ can be described purely by partition combinatorics known as Schubert calculus. There are two important formulas to do computations. First is the complimentary dimension formula which states that if $|\lambda| + |\mu| = r(n-r)$ then $\deg(\sigma_\lambda \sigma_\mu) = 1$ if $\mu = \lambda^c$ and $0$ otherwise. The other formula is called Pieri's formula. For a partition $\lambda \subseteq r\times(n-r)$ and a positive integer $b$ we have $\sigma_\lambda \sigma_{(b)} = \sum_{\mu \in \lambda \otimes b} \sigma_\mu$, where $\lambda \otimes b$ is the set of partitions in $r\times (n-r)$ that can be obtained from the Young diagram of $\lambda$ by adding $b$ boxes, at most one per column.  For a comprehensive reference see \cite{EisHar}.

\subsection{Matroids}
We mainly use the basis definition of a matroid, and identify a matroid $M$ of rank $r$ on $[n]$ with its set of bases $\mathcal{B} \subseteq \binom{[n]}{r}$. For references on matroid theory see \cite{oxleyBook}, \cite{Katz} and \cite{ardila}. We denote the number of connected components of $M$ by $\kappa(M)$.

To a point $x\in G(r,n)$ we associate a matroid $M_x$. The bases of $M_x$ are the sets $I \in \binom{[n]}{r}$ with non-zero Plücker coordinates $p_I(x) \neq 0$. A matroid is called representable over $\CC$ if it is obtained from some point in $G(r,n)$. We will be interested in a subclass of matroids called sparse paving matroids. A matroid $M$ of rank $r$ is called paving if all its circuits have cardinality $r$ or $r+1$. A matroid $M$ is called sparse paving if both $M$ and its dual $M^*$ are paving. It is conjectured in \cite{MAYHEW} and \cite{oxleyBook} that asymptotically almost all matroids are sparse paving.

We will also describe another class of matroids called lattice path matroids in order to make several other definitions easier; see \cite{BoninMierNoy} for comprehensive definitions. Let $P$ and $Q$ be north-east lattice paths from $(0,0)$ to $(n-r,r)$ such that $P$ never goes below $Q$. The lattice path matroid $M[P,Q]$ is a rank $r$ matroid on $[n]$ with bases in bijection with the lattice paths from $(0,0)$ to $(n-r,r)$ bounded by $P$ and $Q$. To get a basis $B$ of $M[P,Q]$ from a lattice path $S$, label the steps of $S$ by the integers $1,\dots, n$ from $(0,0)$ to $(n-r,r)$. Then the basis $B$ is the set of labels of vertical steps.

One special case of lattice path matroids are those where $P$ is the path starting with $r$ vertical steps followed by $n-r$ horizontal steps. In this case the matroids $M[P,Q]$ are called Schubert matroids\footnote{Schubert matroids are known under many names in the literature, such as nested matroids and generalized catalan matroids.}. If $I \in \binom{[n]}{r}$ is the set indicating the vertical steps of $Q$ we denote the Schubert matroid by $SM_I = M[P,Q]$. We introduce the matroids that will be most important to us in the following examples.
\begin{example}
    \label{ex:uniform}
    The Schubert matroid given by $I = \{n-r+1,n-r+2,\dots, n\}$ is the uniform matroid $SM_I = U_{r,n}$. To see this consider the example $r=3$ and $n = 7$. Since the paths $P$ and $Q$ form the $3$ by $4$ rectangle, the bases of $SM_I$ are in bijection with the lattice paths from $(0,0)$ to $(4,3)$. We see that any set $B \in \binom{[7]}{3}$ can be obtained as labels of vertical steps of some path.
\end{example}

\begin{example}
    \label{ex:minimal}
    The Schubert matroid given by $I = \{2,3,\dots,r-1,n \}$ is called the minimal matroid, and is denoted $T_{r,n} = SM_I$. It is, up to isomorphism, the unique connected matroid of rank $r$ on $[n]$ with minimal number of bases. In the example where $r=3$ and $n=7$ it is given by the diagram
    \begin{equation*}
        \scalebox{0.8}{\young(~~~~,~,~)}
    \end{equation*}
    and is represented by the columns of the matrix
    \begin{equation*}
        \left(\begin{array}{ccccccc}
            1 & 0 & 0 & 1 & 1 & 1 & 1\\
            0 & 1 & 0 & 1 & 1 & 1 & 1\\
            0 & 0 & 1 & 1 & 1 & 1 & 1\\
        \end{array}\right).
    \end{equation*}
    In general $T_{r,n}$ is obtained from $U_{r,r+1}$ by adding $n-r-1$ elements parallel to one of the elements in $U_{r,r+1}$.
\end{example}
\begin{example}
    \label{ex:panhandle}
    The Panhandle matroids were introduced in \cite{HMMMNVYPanhandle}. They are given by three positive integers $r \leq s < n$. Here we define them as the Schubert matroid $\Pan_{r,s,n} = SM_I$ where $I = \{s-r+2, s-r+3, \dots ,s, n\}$. The diagram of the panhandle matroid $\Pan_{4,6,10}$ is
    \begin{equation*}
        \scalebox{0.8}{\young(~~~~~~,~~~,~~~,~~~)}
    \end{equation*}
    which looks like a panhandle, hence the name. Notice that $\Pan_{r,r,n} = T_{r,n}$ and $\Pan_{r,n-1,n} = U_{r,n}$.
\end{example}


\subsection{Matroid polytopes and invariants}
In this section we follow the conventions of \cite{DerksenFink} and \cite{FinkSpeyer}. Throughout let $S(r,[n])$ denote the set of matroids of rank $r$ on $[n]$. We are interested in functions $f\colon S(r,[n]) \to A$ with values in an abelian group $A$. Such functions are called matroid invariants if $f(M_1)=f(M_2)$ whenever $M_1$ and $M_2$ are isomorphic. Also $f$ is called additive if $f(M) = 0$ whenever $M$ is not connected. Another important property of functions of matroids is called valuativity, which we will explain in what follows.

To any matroid $M$ on $[n]$ the matroid (base) polytope $P(M) \subset \R^n$ is defined as the convex hull of the indicator vectors $e_B = \sum_{i\in B}e_i$ for all bases $B$ of $M$. In \cite{Gel} it was shown that an equivalent definition of matroids can be given in terms of its matroid polytope, so the matroid polytope $P(M)$ contains all the information of $M$. The dimension of $P(M)$ is $n - \kappa(M)$. We will be interested in how we can chop up matroid polytopes into smaller matroid polytopes.
\begin{definition}
    A matroid subdivision of a matroid $M$ is a polyhedral complex consisting of matroid polytopes with facets $P_1, P_2, \dots P_k$ such that 
    \begin{equation*}
        P(M) = P_1 \cup P_2 \cup \dots \cup P_k.
    \end{equation*}
\end{definition}

Given such a subdivision, for any set $J \subseteq [k]$ we denote $P_J = \bigcap_{j\in J}P_j$ and $P_\emptyset = P(M)$.
We may identify the set $S(r,[n])$ with the set of matroid polytopes of matroids in $S(r,[n])$.

\begin{definition}
    A function $f\colon S(r,[n]) \to A$ is called valuative if for any $M \in S(r,[n])$ and any matroid subdivision of $M$ as above, the relation
    \begin{equation*}
        \sum_{J\subseteq [k]} (-1)^{|J|}f(P_J) = 0
    \end{equation*}
    holds.
\end{definition}
That is, $f$ is valuative if it satisfies an inclusion-exclusion property with respect to matroid subdivisions. Notice that if $f$ is both valuative and additive, since the matroid of $P_J$ is disconnected for all $J$ except $1,2, \dots , k$ and $\emptyset$, the formula simplifies to
\begin{equation*}
    f(M) = \sum_{i=1}^k f(P_i).
\end{equation*}
 In \cite{DerksenFink}, Derksen and Fink proved that a function $f\colon S(r,[n]) \to A$ is valuative if and only if it factors through the abelian group generated by indicator functions of matroid polytopes of $S(r,[n])$, denoted $P(r,[n])$. They also show that the quotient of $P(r,[n])$ given by identifying isomorphic matroids is generated as an abelian group by the classes of the indicator functions of Schubert matroids. This means that any valuative matroid invariant is determined by its values on the Schubert matroids. 

 The most obvious example of a valuative matroid invariant is the volume of a matroid polytope.
 \begin{example}
      Let $M$ be a matroid on $[n]$. We denote the normalized $s$ dimensional volume of $P(M)$ by $\Vol_s(P(M))$. Then
     \begin{equation*}
         \Vol_s(P(-))\colon S(r,[n]) \to \Z
     \end{equation*}
     is a valuative matroid invariant. Moreover, if $s = n-1$ it is an additive matroid invariant since $M$ is connected if and only if $n-\kappa(M) = s$. We may drop the subscript in $\Vol_s(P(M))$ if it does not cause too much confusion.
\end{example}
The next important example was introduced by Crapo in \cite{Crapo}.
\begin{example}
    The beta invariant is a function 
    \begin{equation*}
        \beta \colon S(r,[n]) \to \Z.
    \end{equation*}
    It is completely determined by the properties that if $i$ is neither a loop or coloop of $M$ then
    \begin{equation*}
        \beta(M) = \beta(M/i) + \beta(M\minus i).
    \end{equation*}
    And the base cases: \(\beta(U_{0,1}) = 0\), \(\beta(U_{1,1}) = 1\), and otherwise if $M$ contains a loop or coloop then $\beta(M) = 0$.
    It is well known that $\beta(M)$ is non-negative for all $M$ and $\beta(M) = 0$ if and only if $M$ is disconnected (excluding the special case $U_{0,1}$). The beta invariant is in fact also a valuative and additive matroid invariant.
\end{example}

The next example calculates the beta invariant of minimal matroids.
\begin{example}
    As we have seen, the minimal matroid $T_{r,n}$ is obtained from $U_{r,r+1}$ by adding $n-r-1$ elements parallel to one of the elements in $U_{r,r+1}$. Let $i$ be one of these parallel elements. Then $T_{r,n} \minus i = T_{r,n-1}$ and $T_{r,n}/i$ contains $n-r-1$ loops, so
    \begin{equation*}
        \beta(T_{r,n}) = \beta(T_{r,n}/i) + \beta(T_{r,n} \minus i) = \beta(T_{r,n-1}).
    \end{equation*}
    We may iterate this recursion until $T_{r,m}$ does not contain any parallel elements to choose, i.e.\ when $m-r-1 = 0$. Hence
    \begin{equation*}
        \beta(T_{r,n}) = \beta(T_{r,r+1}) = \beta(U_{r,r+1}).
    \end{equation*}
    The beta invariant of a uniform matroid is known to be
    \begin{equation*}
        \beta(U_{r,n}) = \binom{n-2}{r-1}
    \end{equation*}
    so
    \begin{equation*}
        \beta(T_{r,n}) = \beta(U_{r,r+1}) = \binom{r-1}{r-1} = 1.
    \end{equation*}
\end{example}

The following theorem characterises sparse paving matroids in terms of matroid subdivisions, and can be found for example in \cite{FerroniMinimal} or in \cite{JoswigSchröter}.
\begin{thm}
    A matroid $M \in S(r,[n])$ is sparse paving if and only if there exists a matroid subdivision of $U_{r,n}$ with facets $P(M)$ and $P(M_1), \dots , P(M_k)$ where all matroids $M_i$ are isomorphic to $T_{r,n}$.
\end{thm}
Hence for any additive and valuative matroid invariant $f$ and sparse paving matroid $M$ of rank $r$ on $[n]$ we have
\begin{align*}
    f(U_{r,n}) &= f(M) + k\cdot f(T_{r,n})\\
    f(M) &= f(U_{r,n}) - k\cdot f(T_{r,n}).
\end{align*}

%Notice that we can solve for $k$ using the beta invariant
%\begin{align*}
%    \beta(U_{r,n}) &= \beta(M) + k\cdot \beta(T_{r,n})\\
%    \binom{n-2}{r-1} &= \beta(M) + k\\
%    k &= \binom{n-2}{r-1} - \beta(M).
%\end{align*}